%==============================================================
%  notes/codigo-ejemplos.tex
%  Catálogo visual de las utilidades de estilos/codigo.tex.
%
%  Es un documento INDEPENDIENTE y no forma parte de la tesis: sirve para
%  ver de un golpe todos los estilos disponibles. Compílalo con
%
%      latexmk -pdf codigo-ejemplos.tex
%
%  desde tesis/ o desde tesis/notes/ (ambos funcionan). La referencia
%  escrita está en notes/codigo.md.
%==============================================================
\documentclass[11pt,letterpaper]{book}
\usepackage[utf8]{inputenc}
\usepackage[T1]{fontenc}
\usepackage{lmodern}
\usepackage[spanish,es-tabla]{babel}
\usepackage[margin=2.5cm]{geometry}
\usepackage{color}
\usepackage{hyperref}

% La ruta del estilo cambia según desde dónde se compile.
\IfFileExists{estilos/codigo.tex}
  {\newcommand{\rutacodigo}{estilos/codigo.tex}}
  {\newcommand{\rutacodigo}{../estilos/codigo.tex}}
\input{\rutacodigo}

\begin{document}

\chapter{Catálogo de listados de código}

Cada sección muestra el resultado de una utilidad. El código fuente exacto
de cada ejemplo está en \texttt{notes/codigo.md}.

\section{Python --- \texttt{codigo}, alias \texttt{py}}

\begin{codigo}[py]{Carga del split de prueba de MessIRve}
consulta = "a cuál movimiento literario pertenece fabio fiallo"  # consulta real
datos = load_dataset("spanish-ir/messirve", "full", revision="1.2")["test"]
print(f"{len(datos):,} juicios de relevancia")  # 174,078
\end{codigo}

\section{Shell --- alias \texttt{sh}, y flotante con \texttt{codigoflotante}}

Un listado flotante se etiqueta con la clave \texttt{label} y se referencia como
cualquier figura: el listado~\ref{lst:indexacion}, o con hyperref
\verb|\autoref|: \autoref{lst:indexacion}. (Con cleveref cargado sería
\verb|\cref{lst:indexacion}|.)

\begin{codigoflotante}[sh, label={lst:indexacion}]{Indexación de un modelo base}
python -m ir_spanish.index --modelo bge-m3 --corpus eswiki_20240401_corpus
\end{codigoflotante}

\section{JSON --- alias \texttt{js}}

\begin{codigo}[js]{Configuración del reranker}
{"modelo": "jinaai/jina-reranker-v3", "profundidad": 100, "semilla": 42}
\end{codigo}

\section{Pseudocódigo en español --- alias \texttt{seudo}}

\begin{codigo}[seudo]{Fusión RRF}
algoritmo RRF(listas, k = 60):
    para cada documento d en la unión de las listas:
        puntuacion(d) = suma sobre cada lista L de 1 / (k + rango_L(d))
    devolver los documentos ordenados por puntuacion descendente
\end{codigo}

\section{Salida de consola --- alias \texttt{sal}}

\begin{codigo}[sal]{Evaluación}
$ python -m ir_spanish.evaluate --metrica ndcg@10
nDCG@10 = 0.4123   Recall@100 = 0.8910   MRR@10 = 0.3802
\end{codigo}

\section{Opciones de \texttt{listings} y líneas largas}

Con \texttt{listing options=\{\ldots\}} se pasa cualquier opción de
\texttt{listings}; aquí \texttt{firstnumber=7} y una línea que se parte:

\begin{codigo}[listing options={style=python, firstnumber=7}]{Línea larga}
resultados = evaluar(modelo="multilingual-e5-large-instruct", metrica="ndcg@10", profundidad=100, semilla=42, consulta="ñandú")
\end{codigo}

\section{Inclusión de un archivo real --- \texttt{\textbackslash archivocodigo}}

\begin{verbatim}
\archivocodigo[py, lineas={1}{6}]{../ir-spanish/ruta/al/script.py}{Título}
\end{verbatim}

Aquí se incluyen, como muestra, las primeras líneas de este mismo estilo:

\archivocodigo[lineas={1}{6}]{\rutacodigo}{Encabezado de \texttt{estilos/codigo.tex}}

\section{Código en línea --- \texttt{\textbackslash py}, \texttt{\textbackslash sh}, \texttt{\textbackslash cod}}

Se cargó con \py{revision="1.2"} sobre \sh{ssh iimas} y se evaluó el \cod{top-100}.
La consulta \py{consulta_única} y el término \py{ñandú} conservan sus acentos, igual
que \cod{P@50} y \sh{python -m ir_spanish.evaluate}.

\section{Índice de listados}

Si quieres un «Índice de listados», añade \verb|\lstlistoflistings| en el
material preliminar de \texttt{tesis.tex} (junto a \verb|\listoffigures| y
\verb|\listoftables|); las entradas se registran solas.

\lstlistoflistings

\end{document}
